% --- SLIDE 9: TECHNICAL FEASIBILITY (SINGLE SLIDE) ---
\begin{frame}{Technical Feasibility}
    % Layout pulito con elenchi puntati per massima leggibilità
    \begin{columns}[T]
        \begin{column}{0.95\textwidth}
            \small
            \begin{itemize}
                \item \textbf{1. Kinematic Response (The "Wall Effect" Risk)}
                \begin{itemize}
                    \item \textit{Challenge:} Preventing premature STF hardening during oblique heel strikes at low speeds.
                    \item \textit{Solution:} \alert{Rheological Tuning}. By precisely adjusting the $SiO_2$ mass fraction, we calibrate the \textit{critical shear rate} to ensure activation occurs only during high-energy tasks.
                \end{itemize}
                
                \vspace{0.4cm}
                
                \item \textbf{2. Pendulum Effect (Distal Inertia)}
                \begin{itemize}
                    \item \textit{Challenge:} Fluid mass increases the distal weight, raising metabolic cost during the \textit{Swing Phase}.
                    \item \textit{Solution:} \alert{Spatial Optimization}. STF is not distributed uniformly but encapsulated strictly within primary load-bearing lobes, reducing fluid weight by $\sim$60\% without performance loss.
                \end{itemize}
                
                \vspace{0.4cm}
                
                \item \textbf{3. Material Fatigue \& Leakage}
                \begin{itemize}
                    \item \textit{Challenge:} Cyclic loading ($10^6$ cycles) leading to TPU micro-fractures and fluid leakage.
                    \item \textit{Solution:} \alert{Micro-compartmentalization}. The closed-cell TPMS structure ensures a \textit{fail-safe} design. A localized puncture results in negligible performance degradation instead of catastrophic failure.
                \end{itemize}
            \end{itemize}
        \end{column}
    \end{columns}
\end{frame}